\documentclass[11pt,a4paper]{article}

\usepackage[utf8]{inputenc}
\usepackage[cyr]{aeguill}
\usepackage[francais]{babel}

% Les packages
%\usepackage[french]{babel}
\usepackage{fancybox}
\usepackage{graphics}
\usepackage{color}
\usepackage{latexsym}
\usepackage{amsmath,amsthm,amssymb}
\usepackage[plain]{fullpage}
\usepackage{verbatim}
\usepackage{times,epsfig}
\usepackage{listings}
\usepackage{multirow}
\usepackage{amsmath,amsthm,amssymb,amscd,txfonts,bbm} % RAJOUT ICI !!!
\usepackage{graphics,epsfig} % RAJOUT ICI !!!
%\usepackage{algorithmic}
\usepackage[final]{pdfpages}
\usepackage{fancyvrb}
\usepackage{enumerate}
\usepackage{listingsutf8}
\usepackage{comment}
\usepackage{minted}

% small et verbatim: to be coherent with \file
\newenvironment{smallverbatim}%
{\endgraf\small\verbatim}{\endverbatim}

\newenvironment{qv}
{\quote\Verbatim}
{\endquote\endVerbatim}

\usepackage{lastpage}

\newcommand{\terme}[1]{\texttt{#1}}
\newcommand{\termSet}[1]{\{\texttt{#1}\}}
\newcommand{\guill}[1]{<<\,#1\,>>}
% Macro pour l'inclusion de figure
\newcommand{\fig}[2]{%
  \begin{figure*}[hbt]
   \begin{center}
      \includegraphics[width=12cm]{#1}
  \caption{\label{#1}  #2}
  \end{center}
 \end{figure*}
}
\newcommand{\figTex}[2]{%
  \begin{figure*}[hbt]
 \centerline{\input{#1.pstex_t}}
  \caption{\label{#1}  #2}
 \end{figure*}
}

% Macro résumés
\definecolor{turquoise}{rgb}{0.20 ,0.60, 0.60}
\definecolor{vert}{rgb}{0.40,  0.60 ,0.00}
\definecolor{violet}{rgb}{0.80 , 0.00 , 0.80}
\definecolor{bleu}{rgb}{0.20,  0.20,  0.80}
\definecolor{rose}{rgb}{1.00,  0.20,  0.40}

\lstset{basicstyle=\footnotesize,showstringspaces=false,language=XML,breaklines=false,columns=fullflexible,frame=shadowbox, captionpos=b}

\newenvironment{solution}[0]{
~\\ \color{red} \textbf{Solution : }\color{black}}{%
}
%\excludecomment{solution}


\newcommand{\cle}[1]{\textbf{#1}}
\def\fk#1{\textit{#1}}
\def\tble#1{\textit{#1}}

\parindent=0pt

\usepackage{url}

\newtheorem{Exercice}{Definition}


\begin{document}

\title{Sujet d'examen UE NFE204 - Session 2}
\author{Année universitaire 2021-2022}
\date{Examen session FOAD\,: 1er septembre  2022}

\maketitle

\begin{center}
{\bf \Large Consignes}\\
\begin{tabular}{| p{15cm} |}
\hline
Dur\'ee de l'examen : 3h00; Documents autoris\'es\,: 10 pages recto/verso de
notes manuscrites\,;  Calculatrice autoris\'ee\\

Les exercices sont indépendants les uns des autres. Merci de soigner votre écriture.\\
Ce sujet contient \pageref{LastPage} pages.\\

\textit{Important}\,: nous n'évaluons pas les réponses par la syntaxe. Ne vous
inquiétez pas d'utiliser un point-virgule à la place d'un point d'exclamation ou
autres détails mineurs
\\
\hline
\end{tabular}
\end{center}

%\emph{L'ensemble de l'examen doit vous permettre de montrer que vous avez assimilé les principales notions vues
%en cours et en TP et que vous savez les mettre en application et les expliquer. La clarté, la concision,
%la précision des réponses sera donc appréciée.}


\section*{Première partie : modélisation NoSQL (7 pts)}

Les articles de recherche universitaires sont maintenant gérés dans des dépôts
d'archive, comme HAL en France. Posons-nous quelques questions sur le fonctionnement
d'une telle archive.
Un modèle très sommaire du dépôt est donné dans la figure\,\ref{exam-20-hal}. Les chercheurs
publient des articles (qui peuvent avoir plusieurs co-auteurs) et sont \emph{affiliés} à des organismes
de recherche. Ces organismes sont hiérarchiquement structurés\,: un chercheur peut être
attaché à une équipe (organisme de base), qui fait partie d'un laboratoire (organisme
parent) qui lui-même fait partie d'un établissement (par exemple le Cnam), etc.

\fig{exam-20-hal}{Modèle (sommaire) du dépôt}


\vspace*{0.2cm}

Voici également un tout petit échantillon de la base relationnelle.
\begin{center}
\begin{tabular}{l|l|l}
\cle{ref} &  titre & année\\
\hline
hal-875 & Music modeling & 2018\\
hal-293 & Recommendations on Twitter & 2019\\
\hline
\end{tabular}

\centerline{La table \tble{Article}}


\vspace*{0.2cm}
\begin{minipage}[b]{6cm}
\begin{tabular}{l|l|l}
\cle{id} & nom & affiliation\\
\hline
nt & N. Travers & DVRC \\
pr & P. Rigaux & Vertigo \\
rfs & R. Fournier-S'niehotta & Vertigo \\
cdm & C. du Mouza & ISID \\
\hline
\end{tabular}

\centerline{La table \tble{Chercheur}}
\end{minipage} \hspace*{1cm} \
\begin{minipage}[b]{6cm}
\begin{tabular}{l|l}
\cle{idAuteur} & \cle{refArticle}\\
\hline
hal-875 & nt \\
hal-293 & nt \\
hal-875 & pr \\
hal-875 & rfs \\
hal-293 & cdm \\
\hline
\end{tabular}
\centerline{La table des auteurs}
\end{minipage}
\end{center}

\begin{enumerate}
  \item Proposez un format de document JSON représentant toutes les informations
        relatives à l'article hal-875 présentes dans les trois tables ci-dessus (représentation A).

  \item Proposez un document JSON représentant toutes les
        informations relatives à l'organisme Vertigo présentes dans 3 tables ci-dessus
         (représentation B).

   \item Expliquez le calcul MapReduce permettant de produire les documents
   B à partir des documents A.
   \textbf{Important}\,: pour spécifier les calculs MapReduce, donner les fonctions de Map et de Reduce, en javascript,
en pseudo-code, au pire en langage naturel en étant le plus précis possible.


   \item Maintenant on prend en compte la table décrivant la hiérarchie des organismes,
         dont voici un échantillon.
         \begin{center}
         \begin{tabular}{l|l|l|l}
\cle{id} & intitulé & ville & idParent\\
\hline
Vertigo & Données et apprentissage & Paris & cedric \\
DVRC & Systèmes intelligents & Nanterre & De Vinci \\
ISID & Systèmes d'information & Paris &  cedric\\
cedric & Centre informatique  & Paris &  Cnam\\
\hline
\end{tabular}
\centerline{La table des organismes de recherche}
\end{center}

   Dessinez cette hiérarchie, et proposez une modélisation JSON pour ajouter l'information
   sur les organismes dans la représentation B. Illustrez cette modélisation
   sur le document "Vertigo" de la seconde question.

     \item Les références bibliographiques sont gérées depuis des dizaines d'années dans
       un format texte dit Bibtex. Voici par exemple
       un document Bibtex représentant un article.

\begin{minted}{text}
@article{frt-02435620,
  TITLE = {{Modeling Music as Synchronized Time Series}},
  AUTHOR = {Fournier-S'niehotta, Raphael and Rigaux, Philippe and Travers, Nicolas},
  URL = {https://hal-cnam.archives-ouvertes.fr/hal-02435620},
  JOURNAL = {{Information Systems}},
  PUBLISHER = {{Elsevier}},
  VOLUME = {73},
  PAGES = {35-49},
  YEAR = {2018},
}
\end{minted}

       Comment qualifieriez-vous ce format par rapport à JSON\,? Voyez-vous
        des inconvénients, des avantages, à cette représentation\,?

\end{enumerate}


\section*{Deuxième partie : recherche d'information (8 pts)}

Les articles ont des résumés, et on va construire un index sur ces résumés à des fins
de classification et d'analyse. Voici un extrait des résumés de 5 articles
$A_1, ..., A_5$.
\vspace*{0.3cm}

\begin{tabular}{lp{15cm}}
     $A_1$& Cet article compare les stratégies d'optimisation des échanges réseaux entre processeurs\\
     $A_2$& Un orchestre moderne doit savoir maîtriser une harmonie atonale\\
     $A_3$& Comparaison entre les propriétés des réseaux filaires et des réseaux optiques\\
     $A_4$& Pourquoi un processeur quantique ne peut surmonter un processeur standard
         qu'en présence d'un réseau performant\\
     $A_5$& Mozart est-il mieux servi par un orchestre de chambre ou un orchestre philarmonique?\\
\end{tabular}

\bigskip

On va se limiter au vocabulaire (\texttt{harmonie, processeur, réseau, orchestre}).
(NB: on suppose
     une phase préalable de normalisation qui élimine les pluriels, majuscules, etc.)

\begin{enumerate}
   \item Donnez une matrice d'incidence contenant les tf,
      et un tableau donnant les idf (sans le $\log$), pour les  termes précédents. Placez
      les termes en ligne, les résumés en colonne.

\begin{solution}
\begin{center}
\begin{tabular}{l|c|c|c|c|c}
            & \textbf{A1} & \textbf{A2} & \textbf{A3} &\textbf{A4}&\textbf{A5} \\ \hline
   harmonie 5/1   & 0           & 1           & 0          & 0         & 0\\
   processeur  5/2  & 1           & 0           & 0           & 2        & 0\\
   réseau  5/3  & 1           & 0           & 2          & 1         & 0\\
   orchestre 5/2  & 0           & 1           & 0           & 0        & 2\\\hline
   \end{tabular}
\end{center}
\end{solution}

  \item Donnez les normes des vecteurs représentant ces résumés.
\begin{solution}
Les normes
\begin{itemize}
  \item $||A_1|| = \sqrt{2}$; $||A_2|| = \sqrt{2}$; $||A_3|| = \sqrt{2}$;
    $||A_4|| = \sqrt{5}$, $||A_5|| = \sqrt{4}$
\end{itemize}
\end{solution}

  \item Dans un espace vectoriel à deux dimensions avec "processeur" en abcisse
  et "réseau" en ordonnée, placez les vecteurs représentant nos documents.

  \item Donner les résultats classés par similarité cosinus basée sur les tf
  (on ignore l'idf) pour les requêtes suivantes. Detaillez les calculs.

  \begin{itemize}
    \item processeur et réseau
    \item harmonie et orchestre
  \end{itemize}
\item Sans calcul, indiquez quel document est classé en tête pour la requête
    avec le seul mot "réseau", et expliquez pourquoi.

 \item Indiquez la similarité cosinus de chaque paire de documents.

 \item Des catégories d'article semblent-elles se dégager? Combien, et quelle
 procédure de classification pourriez-vous proposer\,?
\end{enumerate}


\section*{Troisième partie : flux (2 pts)}

Le dépôt reçoit un flux d'articles \texttt{fluxArt} avec leur résumé. On dispose d'un procédure de classification
que nous appellerons \texttt{ClasserArticle(r): C}. Elle prend en entrée un
résumé \texttt{r} et produit une catégorie \texttt{C}.

On veut écrire un traitement de flux qui compte le nombre d'articles
par catégorie. Pour rappel, voici le traitement vu en cours pour compter les mots

\begin{minted}{scala}
case class CompteurMot(mot: String, compteur: Int)
val stream = senv.socketTextStream("localhost", 9000, '\n')
val mots = stream.flatMap ({ str => str.split("\\W+") })
                .map({ CompteurMot(_, 1) })
                .keyBy("mot")
                .reduce( (acc, occ)
                   => {CompteurMot (acc.mot, acc.compteur + 1) })
                .print()
senv.execute("Le compteur de mots")
\end{minted}

Proposez, sur cette base, le traitement comptant le nombre  d'articles
par catégorie. Vous pouvez, si vous le souhaitez, proposer une description
textuelle des  opérateurs successifs à appliquer, ou même un dessin commenté.

\section*{Quatrième partie : questions de cours (3 pts)}

\begin{enumerate}
\item Que signifie, pour une structure de partitionnement, être dynamique.
   Avons-nous  étudié  un système  où  le partitionnement   n'était pas dynamique ?

   \item Quelle sont les opérations proposées par une ressource dans une architecture RESTful ?

    \item Donnez une définition de la scalabilité.
\end{enumerate}

\end{document}
